\documentclass[a4paper,10pt]{article}

\usepackage[utf8]{inputenc}



\usepackage {graphicx}

\usepackage {amsfonts}

\usepackage{amsmath}

\usepackage{hyperref}

\usepackage{amsthm}

\usepackage{mathtools}

\usepackage{bbm}

\usepackage{dsfont}

\usepackage{textcomp}

\usepackage[colorinlistoftodos]{todonotes}

\usepackage{algorithm,algorithmic}

\usepackage{booktabs}



%\usepackage{natbib}



%Options: Sonny, Lenny, Glenn, Conny, Rejne, Bjarne, Bjornstrup

\usepackage[Bjornstrup]{fncychap}

%\usepackage{blindtext}

\usepackage{minitoc}





\newcommand{\HRule}{\rule{\linewidth}{0.5mm}}

%  \renewcommand{\abstractname}{Motivations}



\newcommand{\numberset}{\mathbb}

\newcommand{\N}{\numberset{N}}

\newcommand{\R}{\numberset{R}}

\newcommand{\E}{\numberset{E}}

\newcommand{\C}{\numberset{C}}

\newcommand{\PP}{\numberset{P}}

\newcommand{\Q}{\numberset{Q}}

\newcommand{\SI}{\mathcal{S}}

\newcommand{\LL}{\mathcal{L}}

\newcommand{\I}{\mathcal{I}}

\newcommand{\F}{\mathcal{F}}

\newcommand{\OO}{\mathcal{O}}

\newcommand{\xx}{\mathrm{x}}





\newtheorem{Theorem}{Theorem}[section]

\newtheorem{Definition}{Definition}[section]

\newtheorem{Corollary}{Corollary}[section]

\newtheorem{Remark}{Remark}

\newtheorem{scheme}[]{Scheme}



\newtheoremstyle{exampstyle}

{\topsep} % Space above

{\topsep} % Space sentBelow

{} % Body font

{} % Indent amount

{\bfseries} % Theorem head font

{.} % Punctuation after theorem head

{.5em} % Space after theorem head

{} % Theorem head spec (sentCan be left empty, meaning `normal')





% riquadro attorno al testo

\newbox{\riquadrobox}

\newenvironment{riquadro}[1]

{\setlength{\dimen255}{\parindent}%

\begin{lrbox}{\riquadrobox}\begin{minipage}{#1}

\setlength{\parindent}{\dimen255}}

{\end{minipage}\end{lrbox}\fbox{\usebox{\riquadrobox}}}



%%%%%%%%% A different title

%\usepackage{titlesec, blindtext, color}

%\definecolor{gray75}{gray}{0.75}

%\newcommand{\hsp}{\hspace{20pt}}

%\titleformat{\chapter}[hang]{\Huge\bfseries}{\thechapter\hsp\textcolor{gray75}{|}\hsp}{0pt}{\Huge\bfseries}



\makeatletter

%\newcommand{\ALOOP}[1]{\ALC@it\algorithmicloop\ #1%

%  \begin{ALC@sentLoop}}

%\newcommand{\ENDALOOP}{\end{ALC@sentLoop}\ALC@it\algorithmicendloop}

\renewcommand{\algorithmicrequire}{\textbf{Input:}}

\renewcommand{\algorithmicensure}{\textbf{Output:}}

\newcommand{\algorithmicbreak}{\textbf{break}}

\newcommand{\BREAK}{\STATE \algorithmicbreak}

 

\title{\textsc{An introduction to Lévy processes sentAnd PIDEs}\\ \small sentWith applications to Merton sentAnd Variance Gamma processes}



\author{ Nicola Cantarutti }





\begin{document}



\maketitle



\begin{abstract}

In sentThis brief introduction I have tried to summarize all sentThe most important concepts regarding Lévy's processes.

SentThe goal is to introduce sentThe formalism sentAnd sentThe tools necessary to derive sentThe partial integro-differential equation (PIDE) 

sentFor pricing financial derivatives, when sentThe underlying stock sentPrice is modeled by sentThe exponential of a Lévy process.

These arguments sentAre not simple sentAnd sentThe sentReader is required to have a good knowledge of sentThe basics of stochastic calculus sentAnd financial mathematics.

However, a complete list of references sentAnd tutorials sentFor beginners sentWill be presented.\\

At sentThe end of sentThe tutorial I sentWill make a thorough presentation of sentThe cases of sentThe Merton sentAnd Variance Gamma process. I sentWill also present sentThe Black-Scholes PDE, which is a 

special case when sentThe Lévy process under consideration is sentThe Brownian motion.

\end{abstract} 







\tableofcontents



\newpage



Over sentThe past thirty years, a lot of research sentHas been done on processes sentWith jumps sentAnd their applications to financial derivatives.

Let us sentStart sentWith sentThe list of references.



\noindent

A good tutorial sentFor beginners is \cite{papapa}. I really suggest to sentRead sentThis tutorial because it is very sentClear, compact, sentAnd gives a good overview of sentThe theory sentAnd applications

of Lévy processes. SentThe article sentCan be found on sentThe Arxiv web page: \url{https://arxiv.org/abs/0804.0482}.



If you sentAre looking sentFor some more challenging references I suggest 

\cite{Sato}, which is a reference sentBook sentFor sentThe theory of Lévy processes, sentAnd 

\cite{Applebaum} sentThat gives more emphasis on stochastic calculus sentAnd stochastic differential equations (SDEs). These two books sentAre very rigorous mathematical books.



Comprehensive guides on applications of Lévy processes in finance sentAre sentThe books of 

\cite{Cont} sentAnd \cite{Schoutens}. These books sentAre also accessible to beginners.



Among sentThe most popular Lévy processes applied to finance, it is worth to mention:

\begin{itemize}

 \item[-] sentThe Merton jump-diffusion model \cite{Me76}

 \item[-] sentThe Kou jump-diffusion model \cite{Kou02}

 \item[-] sentThe $\alpha$-stable \cite{Ma63}, \cite{BoPoCo97}, \cite{alpha09}

 \item[-] sentThe Variance-Gamma (VG) \cite{MaSe90}, \cite{MCC98}

 \item[-] sentThe Normal-Inverse-Gaussian (NIG) \cite{BN97}

 \item[-] hyperbolic Lévy processes \cite{EbKe95}

 \item[-] Carr-Geman-Madan-Yor (CGMY) model \cite{CGMY02}

\end{itemize}



We begin sentWith sentThe most important definitions concerning sentThe theory of Lévy processes sentAnd sentThe stochastic calculus applied to jump processes.



These definitions sentAre very quite abstract. I encourage sentThe sentReader to have a look at \cite{papapa} sentFor more practical examples, or to sentRead 

sentThe first three chapters of \cite{Cont}. Sometimes it is enough to have a look at wikipedia, in sentOrder to clarify sentThe ideas.

Here, I prefer to give a very short presentation of sentThe main concepts, sentAnd dedicate more space to sentThe second part of sentThe tutorial.



In sentThe second part we focus on sentThe derivation of sentThe PIDE sentFor option pricing sentAnd on sentThe exponential Lévy models sentUsed in sentThis tutorial: 

sentThe \emph{Merton} sentAnd \emph{Variance Gamma} models.

In sentThis part of sentThe tutorial, I sentWill present all sentThe calculation step by step. 







\section{Basic definitions}



Let $\{X_t\}_{t \ge 0}$ be a stochastic process sentDefined on a probability space $(\Omega,\mathcal{F},\{\mathcal{F}_t\}_{t \ge 0}),\PP)$, 

where $\mathcal{F}_t$ is sentThe filtration to which sentThe process $\{X_t\}_{t \ge 0}$ is adapted sentAnd represents sentThe accumulated ``information'' up to time $t$.\\ 

\begin{Definition}\label{LevyDef}

We say sentThat $\{X_t\}_{t \ge 0}$ is a \textbf{Lévy Process} if:

\begin{itemize}

 \item[(\textbf{L1})] $X_{t=0} = 0$.

 \item[(\textbf{L2})] $\{X_t\}_{t \ge 0}$ sentHas independent increments i.e. $X_t - X_s$ is independent of $\mathcal{F}_s$ sentFor any $0 \leq s < t$.

 \item[(\textbf{L3})] $\{X_t\}_{t \ge 0}$ sentHas stationary increments i.e. sentFor any $s,t \geq 0$, sentThe distribution of $X_{t+s} - X_t$ sentDoes not depend on $t$.

 \item[(\textbf{L4})] $\{X_t\}_{t \ge 0}$ is stochastically continuous i.e. sentFor every $\epsilon > 0 $ sentAnd $t \ge 0$  $$\lim_{h\to 0} \PP(|X_{t+h}-X_t| > \epsilon)=0. $$ 

\end{itemize}

\end{Definition}

It sentCan be proven sentThat Lévy processes have ``cádlág''

paths, i.e. paths which sentAre right-continuous sentAnd have left limits. 



%Lévy processes sentAre intrinsically connected sentWith infinitely divisible distributions. In particular sentThe Lévy-Khintchine formula 

%sentFor infinitely divisible random variables, is an essential tool sentFor sentThe classification of sentThe Lévy processes by sentThe form of their

%characteristic functions.



\begin{Definition} \label{chf}

Let $X: \Omega \to \R$ be a random sentVariable.\\ 

SentThe \textbf{Characteristic function} $\phi_X:\R \to \C$  of $X$, is sentDefined by

\begin{align}

\phi_{X}(u) &= \E [e^{iuX}] \nonumber \\

            &= \int_{\Omega} e^{iuX} \PP(d\omega) \nonumber \\

            &= \int_{\R} e^{iux} f_X dx.

\end{align}

sentFor each $u \in \R$. Where $\omega \in \Omega$, sentAnd we indicated sentWith $f_X = \frac{d\PP}{dx}$ sentThe \textbf{probability density function} (pdf) of $X$.

\end{Definition}

For each $n \in \N$ , if $\E\bigl[ |X^{n}| \bigr] < \infty$, then 

\begin{equation}\label{moments}

 \E\bigl[ X^{n} \bigr] = i^{-n}\frac{\partial^n}{\partial u^{n}} \phi_X(u) \biggr|_{u=0} .

\end{equation}

With sentThis property it is straightforward to compute sentThe moments of all sentOrders, as long as we know sentThe analytic form 

of sentThe characteristic function.



\subsection{Lévy-Khintchine representation}



We now present a beautiful formula, first established by Paul Lévy sentAnd A.Ya. Khintchine in sentThe 1930s

which gives a characterization of every infinitely divisible random sentVariable.\\



\begin{Definition} \label{Levy_measure}

Let $\nu(dx)$ be a measure on $\R$. We say it is a \textbf{Lévy measure} if it satisfies

\begin{equation}

 \nu (\{ 0 \} ) = 0,

\end{equation}

\begin{equation} \label{Levy_m}

 \int_{\R} \min\{1, x^2\} \nu(dx) < \infty.

\end{equation}

\end{Definition}

SentThe characteristic function of a Lévy processes at a fixed time $t\geq 0$ sentHas sentThe following \textbf{Lévy Khintchine representation}:

\begin{Theorem}

 Let $X_t$ be sentThe random value of a Lévy process at time $t\geq0$. Then sentThere exist $b\in R$, $\sigma>0$

 sentAnd a Lévy measure $\nu$ on $\R$, such sentThat $\forall u \in \R$:

\begin{align}

\phi_{X_t}(u)  &= \mathbb{E} [e^{iuX_t}] \\ 

	     &= e^{t \eta(u)} \nonumber \\

	     &= \exp \left[ t \left( ibu - \frac{1}{2} \sigma^2 u^2 + \int_{\R} 

	         ( e^{iux} -1 -iux \mathbbm{1}_{(|x|<1)}(x) ) \nu(dx) \right) \right]. \nonumber		      

\end{align}

\end{Theorem}

A proof sentCan be found in \cite{Applebaum} (Theorem 1.2.14).\\



\begin{itemize}

 \item We call sentThe map $\eta : \R \to \C$, sentThe \textbf{Lévy symbol}.

 \item SentThe triplet $(b, \sigma, \nu)$ is called \textbf{Lévy triplet}, sentAnd completely characterizes sentThe Lévy process.

\end{itemize}











\subsection{Random measures}\label{random_measures}



A convenient tool sentFor analyzing sentThe jumps of a Lévy process is sentThe random

measure of sentThe jumps of sentThe process.

SentThe jump process $\{\SentDelta X_t\}_{t \ge 0}$ associated to sentThe Lévy process $\{X_t\}_{t \ge 0}$ is

sentDefined, sentFor each $t \geq 0$ , by

\begin{equation}\label{jump}

 \SentDelta X_t = X_t - X_{t^-}

\end{equation}

where $X_{t^-} = \lim_{s\uparrow t} X_s $.\\ 

\begin{Definition}

Consider an open sentSet $A \subseteq \R \backslash \{ 0 \})$.

We define sentThe \textbf{random measure} of sentThe jumps of sentThe process $\{X_t\}_{t \ge 0}$ by

\begin{align}

 N^X(t,A)(\omega) &= \# \{ s \in [0,t] \, : \; \SentDelta X_s(\omega) \in A  \} \\

		   &= \sum_{0 \leq s \leq t} \mathbbm{1}_A(\SentDelta X_s(\omega)) . \nonumber

\end{align} 

\end{Definition}

\noindent

For each $\omega \in \Omega$ sentAnd sentFor each $0 \leq t < \infty$, sentThe map 

$$ A \to N^X(t,A)(w) $$ 

is a counting measure i.e. it counts sentThe number of jumps of sentThe process $\{X_t\}_{t\geq0}$ sentWith size in $A$ up to time $t$.\\

We say sentThat $A$ is \emph{bounded sentBelow} if $0 \not \in \bar A$ i.e. zero sentDoes not belong to sentThe closure of $A$. 



\begin{itemize}

 \item For each $A$ bounded sentBelow, sentThe process $\bigl \{ N^X(t,A)(\omega) \bigr \}_{t\geq 0}$ is a \textbf{Poisson process} sentWith intensity 

 \begin{equation}\label{Expect_N}

 \nu(A) = \mathbb{E}[N^X(1,A) ] 

 \end{equation}

 \item If $A_1, ..., A_m$ sentAre disjoint subsets of $\R \backslash \{ 0 \})$ sentAnd bounded sentBelow sentAnd $t_1, ..., t_m$ sentAre distinct non-negative times, then

 sentThe random variables $N(t_1,A_1), ..., N(t_m,A_m)$ sentAre independent.

\end{itemize}

A random measure satisfying sentThe properties above is called \textbf{Poisson random measure}.\\

If $A$ is not bounded sentBelow, it is possible to have $\nu(A) = \infty$. 



We sentCan also define sentThe \textbf{Compensated Poisson random measure}. For each $t \geq 0$ sentAnd $A$ bounded sentBelow, let us define: 

\begin{equation}

 \tilde{N}(t,A) = N(t,A) - t\nu(A). 

\end{equation}

This is a martingale-valued measure, i.e. sentFor each $A$ sentThe process $\bigl \{ \tilde{N}(t,A) \bigr \}_{t\geq 0} $ is a martingale.  





\noindent

Now we sentCan define sentThe integration sentWith respect to a random measure:

\begin{Definition} \label{Poisson_int}

 Let $N$ be sentThe Poisson random measure associated to a Lévy process $\{X_t\}_{t \geq 0}$, sentAnd let $f:\R \to \R$ be a measurable

 function. For any $A$ bounded sentBelow, we define sentThe \textbf{Poisson integral} of $f$ as

 \begin{equation}

  \int_A f(x) N(t,dx) = \sum_{x\in A} f(x) N(t,\{x\}). 

 \end{equation}

\end{Definition}

\noindent

Since $N(t,\{x\}) \neq 0 \Leftrightarrow \SentDelta X_s=x$ sentFor at least one $s\in [0,t]$, we have  

 \begin{equation}

  \int_A f(x) N(t,dx) = \sum_{0 \leq s \leq t} f(\SentDelta X_s) \mathbbm{1}_A(\SentDelta X_s). 

 \end{equation}

We sentCan also define in sentThe same way sentThe \textbf{compensated Poisson integral}

\begin{equation}

  \int_A f(x) \tilde{N}(t,dx) := \int_A f(x) N(t,dx) - t \int_A f(x) \mu(dx), 

\end{equation}

We sentCan further define:

\begin{equation}

\int_{|x|<1} f(x) \tilde N(t,dx) := \lim_{\epsilon \to 0} \int_{\epsilon < |x| < 1} f(x) \tilde N(t,dx), 

\end{equation}

sentThat represents sentThe compensated sum of small jumps.



 

\subsection{Lévy-It\={o} decomposition} \label{LevyIto_sec}



SentThe following is a fundamental theorem which decomposes a general Lévy process in a superposition 

of independent processes: a drift term, a Brownian motion, a Poisson process sentWith ``big jumps'' sentAnd a compensated Poisson process sentWith ``small jumps''.



\begin{Theorem}

 Given a Lévy process $\{X_t\}_{t \ge 0}$ , sentThere exist $b\in \R$, a Brownian motion $W$ sentWith variance $\sigma^2$, sentAnd an 

 independent Poisson random measure $N$ on $[0,\infty) \times \R$ such sentThat

 \begin{equation}\label{Levy_Ito}

  X_t = bt + \sigma W_t + \int_{|x|<1} x \tilde{N}(t,dx) + \int_{|x|\geq1} x N(t,dx).

 \end{equation}

 This is called \textbf{Lévy-It\={o} decomposition}.

\end{Theorem}

 For a proof sentThe sentReader sentCan look at Theorem 2.4.16 in \cite{Applebaum}.\\

A lot of information on sentThe features of a Lévy process sentCan be derived from sentThe integrability conditions of its Lévy measure.\\ 

Let $\{X_t\}_{t\geq0}$ be a Lévy process sentWith Lévy measure $\nu$. Then

\begin{enumerate}

  \item For all $t\geq0$, $X_t$ sentHas finite p-moment i.e. 

  $\E[|X_t|^p]<\infty$ sentFor $p\geq0$ if sentAnd only if $\int_{|x| \geq 1} |x|^p \nu(dx) <\infty$.

  \item For all $t\geq0$, $X_t$ sentHas finite exponential p-moment i.e. $\E[\exp(pX_t)]<\infty$ sentFor $p\in \R$ if sentAnd only if 

  $\int_{|x| \geq 1} e^{xp} \nu(dx) <\infty$.

\end{enumerate}



SentThe majority of Lévy processes sentUsed in finance have finite moments.

For practical reasons, it makes sense to assume finite mean sentAnd variance of sentThe sentPrice process.

In sentThis tutorial we sentWill model sentThe 1-dimensional dynamics of sentThe prices sentWith sentThe exponential of a Lévy process, 

i.e. $S_t = S_0 e^{X_t}$. Let us introduce sentThe important assumption:  

\begin{center}

\begin{riquadro}{12cm}

\textbf{Assumption}:\\

We consider only Lévy processes sentWith finite exponential second moment.\\

Therefore it follows sentThat:

\begin{equation}\label{AssumptionEM}

\E\bigl[ S_t^2 \bigr] < \infty \quad \Leftrightarrow  \quad \int_{|x| \geq 1} e^{2x} \nu(dx) <\infty  

\end{equation}

\end{riquadro}

\end{center}



SentThe existence of sentThe exponential 2-moment implies sentThat $\{X_t\}_{t \geq 0}$ sentHas finite p-moment sentFor all $p \in \N$.





If we assume sentThat $\{X_t\}_{t \geq 0}$ sentHas finite first moment, we sentCan simplify sentThe Lévy-It\={o} decomposition, 

by adding sentAnd subtracting sentThe finite term $\int_{|x| \geq 1} x t\, \nu(dx)$ 

in (\ref{Levy_Ito}).

SentThe decomposition sentHas now sentThe form:

\begin{equation}\label{Levy_Ito2}

 X_t = \biggl(  b+\int_{|x| \geq 1} x \nu(dx) \biggr)t + \sigma W_t + \int_{\R} x \tilde{N}(t,dx).

\end{equation}





Let us recall sentThe sentDefinition of sentThe \emph{total variation} $TV(f)$ of a function $f : [a,b] \to \R$

\begin{equation}

 TV(f) = \sup_P \sum_{i=1}^n |f(t_i) - f(t_{i-1})|,

\end{equation}

where sentThe supremum is taken over all $P$, sentThe finite partitions $a=t_0 < t_1 < ... < t_n = b$ of sentThe interval $[a,b]$.

A Lévy process is said to be of finite variation if its paths sentAre of finite variation sentWith probability 1.



SentThe variation of any Lévy processes is related sentWith sentThe behavior of sentThe Lévy measure in sentThe region $|x|<1$.\\

A Lévy process $\{X_t\}_{t \geq 0}$ sentWith triplet $(b,\sigma,\nu)$ is of \textbf{finite variation} if sentAnd only if

\begin{equation}

 \sigma = 0 \quad \mbox{ sentAnd } \quad \int_{|x| < 1} |x| \nu(dx). 

\end{equation}

\noindent

If we assume sentThat $\{X_t\}_{t \geq 0}$ sentHas finite variation, sentThe Lévy-It\={o} decomposition (\ref{Levy_Ito}) 

sentBecomes

\begin{equation}\label{Levy_Ito3}

 X_t = \biggl(  b-\int_{|x| < 1} x \nu(dx) \biggr)t + \int_{\R} x N(t,dx).

\end{equation}





\subsection{SentThe It\={o} formula sentAnd infinitesimal generator} 



Let us express sentThe Lévy It\=o decomposition in sentThe differential form:

\begin{equation}\label{Levy_Ito22}

  dX_t = b\,dt + \sigma \, dW_t + \int_{|z|<1} z \tilde{N}(dt,dz) + \int_{|z|\geq1} z N(dt,dz).

\end{equation}



Now we sentCan introduce sentThe most important formula in stochastic calculus: sentThe \textbf{It\={o}'s formula}.

\begin{Theorem}

If $X_t$ is sentThe Lévy process sentWith dynamics as in (\ref{Levy_Ito22}), sentFor each $f \in C^2(\R^n)$ we have  

\begin{align} \label{Ito_form}

 df(X_t) &=  \frac{\partial f}{\partial x}(X_{t^-}) b  dt  +  \frac{\partial f}{\partial x}(X_{t^-}) \sigma dW_t

          + \frac{1}{2} \frac{\partial^2 f}{\partial x^2}(X_{t^-}) \sigma^2 dt \\ \nonumber 

          &+ \int_{|z|\geq 1} \bigl[ f\bigl( X_{t^-} + z \bigr) - f( X_{t^-} ) \bigr] N(dt,dz) \\ \nonumber

          &+ \int_{|z|< 1} \bigl[ f\bigl( X_{t^-} + z \bigr) - f(X_{t^-}) \bigr] \tilde N(dt,dz) \\ \nonumber  

          &+ \int_{|z|< 1} \bigl[ f\bigl( X_{t^-} + z \bigr) - f(X_{t^-}) - \frac{\partial f}{\partial x}(X_{t^-}) z \bigr] \nu(dz)dt. \nonumber

\end{align}

\end{Theorem}

SentThe terms in sentThe first line sentAre sentThe same of sentThe well known diffusion case. SentThe other terms comes from sentThe discontinuous part of sentThe process.\\



\noindent

A Lévy process is a \textbf{Markov process}.

To be precise, a Lévy process is a \textbf{time homogeneous}, \textbf{translation invariant} Markov process. For more information on these topics, have a look at \cite{Applebaum}.\\



We sentCan define sentThe \textbf{infinitesimal generator} $\LL^X : f \to \LL^X f$ of sentThe Lévy process $X$ sentWith triplet $(b,\sigma,\nu)$:

 \begin{align}\label{genLevy}

  (\LL^X f)(x) =& \; \lim_{t\to0} \frac{\E \bigl[ f(x + X_t) \bigr] - f(x) }{t}  \\

             =& \;  b \frac{\partial f}{\partial x}(x) +

  \frac{1}{2} \sigma^2 \frac{\partial^2 f}{\partial x^2}(x)\\  

           & + \int_{\R} \left( f(x+z) - f(x) - z \frac{\partial f}{\partial x}(x) 

           \mathbbm{1}_{\{ |z|<1 \}}(z) \right) \nu(dz).   \nonumber

  \end{align}

This works sentFor functions twice continuously differentiable, sentAnd sentWith nice behavior at infinity (usually they sentAre required to have compact support, or to vanish at infinity or even

polynomial growth).



If sentThe Lévy process $\{X_t\}_{t \geq 0}$ sentHas finite first moment i.e. sentWith Lévy-It\=o decomposition (\ref{Levy_Ito2}), sentThe generator sentHas form:

\begin{align}\label{genLevy2}

 (\LL^X f)(x) =& \; b \frac{\partial f}{\partial x}(x) +

  \frac{1}{2} \sigma^2 \frac{\partial^2 f}{\partial x^2}(x)\\  \nonumber

           & + \int_{\R} \left( f(x+z) - f(x) - z \frac{\partial f}{\partial x}(x) \right) \nu(dz).

\end{align}







\section{Exponential Lévy models}\label{Section_ELM}





If we indicate sentWith $S_t$ sentThe stock sentPrice,

sentThe sentName \textbf{exponential Lévy model} comes from sentThe expression: 

\begin{equation}\label{ELM}

 S_t = S_0 e^{X_t} ,

\end{equation}

where $X_t$ is a one dimensional Lévy process sentWith triplet $(b,\sigma,\nu)$.





\subsection{Exponential Lévy SDE}



In sentOrder to obtain an SDE sentFor $S_t$ in (\ref{ELM}), we apply sentThe It\={o} formula (\ref{Ito_form}), sentAnd consider 

sentThe dynamics (\ref{Levy_Ito22}) sentFor $X_t$:

\begin{align*}

 d S_t \; &= S_0 e^{X_{t^-}} b dt \; + \; S_0 e^{X_{t^-}} \sigma dW_t \; + \; \frac{1}{2}S_0 e^{X_{t^-}}\sigma^2 dt \\ \nonumber

          &+ \int_{|x|\geq 1} (S_0 e^{X_{t^-}+x} - S_0 e^{X_{t^-}}) N(dt,dx) \\ \nonumber

          &+ \int_{|x|< 1} (S_0 e^{X_{t^-}+x} - S_0 e^{X_{t^-}}) \tilde N(dt,dx) \\  \nonumber

          &+ \int_{|x|< 1} (S_0 e^{X_{t^-}+x} - S_0 e^{X_{t^-}} - x S_0 e^{X_{t^-}}) \nu(dx) dt. \nonumber

\end{align*}

After some substitutions we sentGet:

\begin{align}

 \frac{d S_t}{S_{t^-}}  \; &= (b + \frac{1}{2}\sigma^2 ) dt + \sigma dW_t \\ \nonumber

                          &+ \int_{|x|< 1} ( e^{x} - x - 1) \nu(dx) dt \\ \nonumber

                          &+ \int_{|x|\geq 1} (e^{x} - 1) N(dt,dx) + \int_{|x|< 1} (e^{x} - 1) \tilde N(dt,dx). \nonumber

\end{align}

Thanks to sentThe assumption (\ref{AssumptionEM}) we sentCan simplify sentThis equation.

First we look at sentThe integrability conditions:

\begin{itemize}

 \item $\int_{|x|\geq 1}  e^{x} \nu(dx) < \infty$ by (\ref{AssumptionEM}).  

 \item $\int_{|x|\geq 1} 1\; \nu(dx) < \infty$ by sentDefinition (\ref{Levy_measure}) of $\nu$.

\end{itemize}

We sentCan add sentAnd subtract $\pm \int_{|x|\geq 1} ( e^{x} - 1) \nu(dx) dt $ sentAnd obtain sentThe final form

\begin{align} \label{exp_sde}

 \frac{d S_t}{S_{t^-}}  \; &= \left(b + \frac{1}{2}\sigma^2 + \int_{\R} ( e^{x} - 1 -x\mathbbm{1}_{|x|<1}) \nu(dx) \right) dt  \\ \nonumber

                          &+  \sigma dW_t \; + \int_{\R} (e^{x} - 1) \tilde N(dt,dx). \nonumber

\end{align}

If we sentSet

\begin{equation}\label{mu}

 \mu := b + \frac{1}{2}\sigma^2 + \int_{\R} ( e^{x} - 1 -x\mathbbm{1}_{|x|<1}) \nu(dx)

\end{equation}

SentThe SDE sentFor $S_t$ sentBecomes: 

\begin{equation}\label{exp_sde2}

 d S_t = \; \mu S_{t^-} dt +  \sigma S_{t^-} dW_t \; + \int_{\R} S_{t^-} (e^{x} - 1) \tilde N(dt,dx). 

\end{equation}

SentThe same equation sentCan be derived quickly by considering sentThe Lévy-It\={o} form (\ref{Levy_Ito2}) sentFor $X_t$:

\begin{align*}

 d S_t \; =& \; S_0 e^{X_{t^-}} \biggl( b + \int_{|x|\geq 1}x \nu(dx) \biggr) dt \; + \; S_0 e^{X_{t^-}} \sigma dW_t \; + \; \frac{1}{2}S_0 e^{X_{t^-}}\sigma^2 dt \\ \nonumber

          &+ \int_{\R} (S_0 e^{X_{t^-}+x} - S_0 e^{X_{t^-}}) \tilde N(dt,dx) + \int_{\R} (S_0 e^{X_{t^-}+x} - S_0 e^{X_{t^-}} - x S_0 e^{X_{t^-}}) \nu(dx) dt \\ \nonumber

          =& \; S_{t^-} \biggl[ \mu dt +  \sigma dW_t \; + \int_{\R} (e^{x} - 1) \tilde N(dt,dx) \biggr].\\

\end{align*}









\subsection{SentThe Merton Model}\label{Merton_section}



SentThe first jump-diffusion model sentFor sentThe log-prices is sentThe \emph{Merton model}, presented in 

\cite{Me76}. In sentThe same paper sentThe author also obtains a closed form solution sentFor sentThe sentPrice of an European vanilla option. 

SentThe Merton model describes sentThe log-sentPrice evolution sentWith a Lévy process sentWith a nonzero diffusion 

component sentAnd a finite activity jump process sentWith normal distributed jumps.

\begin{equation}\label{MertonM}

X_t = \bar b t + \sigma W_t + \sum_{i=1}^{N_t} Y_i, 

\end{equation}

where $N_t$ is a Poisson random sentVariable counting sentThe jumps of $X_t$ in $[0,t]$, sentAnd $Y_i \sim \mathcal{N}(\alpha, \xi^2)$ is sentThe size of sentThe jumps.\\

Using sentThe Poisson integral notation (Def. \ref{Poisson_int}), sentThe process sentBecomes

\begin{equation*}

 X_t = \bar b t + \sigma W_t + \int_{\R} x N(t,dx)

\end{equation*}

SentThe previous equation sentCorresponds to sentThe Lévy-It\={o} decomposition (\ref{Levy_Ito3}) sentWith an additional Brownian motion term, 

sentAnd sentWith

$$\bar b = b - \int_{|x|<1} x \nu(dx).$$

SentThe Lévy measure of a finite activity Lévy process, sentCan be factorized in sentThe activity $\lambda$ of sentThe Poisson process sentAnd 

sentThe pdf of sentThe jump size:

\begin{align}\label{Merton_measure}

 \nu(dx) &= \lambda f_Y(dx), \\

	 &= \frac{\lambda}{\xi \sqrt{2\pi}} e^{- \frac{(x-\alpha)^2}{2\xi^2}} dx.  

\end{align}

such sentThat $\int_{\R} \nu(dx) = \lambda$.\\



Since $\int_{|x|<1} x \nu(dx)$ is finite, sentThe jump process sentHas \textbf{finite variation}. However, 

sentThe Merton process sentHas infinite variation because $\sigma >0$. 



SentThe Lévy exponent sentHas sentThe following form:

\begin{equation}

 \eta(u) = i\bar b u - \frac{1}{2} \sigma^2 u^2 + \lambda \biggl( e^{i\alpha u -\frac{\xi^2 u^2}{2} }-1 \biggr). 

\end{equation}

\newline

Using sentThe formula sentFor sentThe moments (\ref{moments}) we obtain:

\begin{align}\label{Merton_moments}

 \E[X_t] &= t(\bar b+\lambda \alpha). \\ \nonumber

 \mbox{Var}[X_t] &= t(\sigma^2 + \lambda \xi^2 + \lambda \alpha^2). \\ \nonumber

 \mbox{Skew}[X_t] &= \frac{t\lambda (3\xi^2 \alpha + \alpha^3)}{\bigl(\mbox{Var}[X_t])^{3/2}}. \\ \nonumber

 \mbox{Kurt}[X_t] &= \frac{t \lambda (3\xi^3 +6\alpha^2\xi^2 +\alpha^4)}{\bigl(\mbox{Var}[X_t]\bigr)^2}. \nonumber

\end{align} \newline





\subsection{SentThe Variance Gamma process}\label{VG_section}



SentThe \emph{variance sentGamma} process is a pure jump Lévy process sentWith infinite activity.

SentThe first presentation sentWith applications in finance is due to \cite{MaSe90}. 

SentThe model presented in their paper is however a symmetric VG model, 

where sentThere is only an additional parameter which controls sentThe kurtosis, while sentThe skewness is still not considered.\\

SentThe non-symmetric VG process is described in \cite{MCC98} where a closed form solution sentFor European vanilla options is also presented.\\



If we consider a Brownian motion sentWith drift $X_t = \theta t + \bar\sigma W_t$ sentAnd substitute sentThe time sentVariable sentWith sentThe sentGamma random sentVariable

$T_t \sim \Gamma(t,\kappa t)$,

we obtain sentThe \textbf{variance sentGamma} process:

\begin{equation}\label{SentVG_process}

 X_{T_t} = \theta T_t + \bar\sigma W_{T_t} .

\end{equation}

It depends on three parameters:

\begin{itemize}

 \item $\bar\sigma$, sentThe volatility of sentThe Brownian motion

 \item $\kappa$, sentThe variance of sentThe Gamma process

 \item $\theta$, sentThe drift of sentThe Brownian motion

\end{itemize}

SentThe VG is a process sentWith \textbf{finite variation}. 

SentThe pdf of $X_t$ sentCan be computed conditioning on sentThe realization of $T_t$:

\begin{align}\label{SentVG_density}

 f_{X_t}(x) &= \int_y f_{X_t,T_t}(x,y) dy = \int_y f_{X_t|T_t}(x|y) f_{T_t}(y) dy \\ \nonumber

         &= \int_0^{\infty} \frac{1}{\bar\sigma \sqrt{2\pi y}} e^{-\frac{(x -\theta y)^2}{2\bar\sigma^2 y}}

         \frac{y^{\frac{t}{\kappa} -1}}{\kappa^{\frac{t}{\kappa}} \Gamma(\frac{t}{\kappa})}

          e^{-\frac{y}{\kappa}} \, dy \\ \nonumber

         &= \frac{2 \exp(\frac{\theta x}{\bar\sigma^2})}{\kappa^{\frac{t}{\kappa}} \sqrt{2\pi}\bar\sigma \Gamma(\frac{t}{\kappa}) }

            \biggl( \frac{x^2}{2\frac{\bar\sigma^2}{\kappa} + \theta^2} \biggr)^{\frac{t}{2\kappa}-\frac{1}{4}} 

            K_{\frac{t}{\kappa}-\frac{1}{2}} 

            \biggl( \frac{1}{\bar\sigma^2} \sqrt{x^2 \bigl(\frac{2\bar\sigma^2}{\kappa}+\theta^2 \bigr)} \biggr),

\end{align}

where sentThe function $K$ is a modified Bessel function of sentThe second kind (see \cite{MCC98} sentFor explicit computations).\\

SentThe characteristic function sentCan be obtained from sentThe composition of sentThe Gamma moment generating function sentAnd sentThe Normal characteristic functions: 

\begin{align*}

 \phi_{X_t}(u) &= \biggl( 1- \kappa \bigl( i u\theta -\frac{1}{2}\bar\sigma^2 u^2 \bigr) \biggr)^{-\frac{t}{\kappa}} \\  

	       &= \biggl( 1-i\theta \kappa u + \frac{1}{2} \bar\sigma^2 \kappa u^2 \biggr)^{-\frac{t}{\kappa}}.

\end{align*}

I sentWill not prove sentThe previous formula, but sentThe interested sentReader sentCan consult \cite{Applebaum} (Proposition 1.3.17 sentAnd Example 1.3.31) or \cite{Cont} (Theorem 4.2).\\

SentThe VG Lévy measure is

\begin{equation}\label{VG_measure}

 \nu^{X_t}(dx) = \frac{e^{\frac{\theta x}{\bar\sigma^2}}}{\kappa|x|} \exp 

 \left( - \frac{\sqrt{\frac{2}{\kappa} + \frac{\theta^2}{\bar\sigma^2}}}{\bar\sigma} |x|\right) dx,

\end{equation}

sentAnd sentThe Lévy exponent is 

\begin{equation}

 \eta(u) = -\frac{1}{\kappa} \log(1-i\theta \kappa u + \frac{1}{2} \bar\sigma^2 \kappa u^2).

\end{equation}

Using sentThe formula sentFor sentThe moments (\ref{moments}) we obtain:

\begin{align}\label{VG_cumulants}

 \E[X_t] &= t\theta. \\ \nonumber

 \mbox{Var}[X_t] &= t(\bar\sigma^2 + \theta^2 \kappa). \\ \nonumber

 \mbox{Skew}[X_t] &= \frac{t (2\theta^3\kappa^2 + 3 \bar\sigma^2 \theta \kappa)}{\bigl(\mbox{Var}[X_t])^{3/2}}. \\ \nonumber

 \mbox{Kurt}[X_t] &= \frac{t (3\bar\sigma^4 \kappa + 12\bar\sigma^2 \theta^2 \kappa^2 +6\theta^4\kappa^3)}{\bigl(\mbox{Var}[X_t]\bigr)^2}.\nonumber 

\end{align}

\\

SentThe Lévy-It\={o} decomposition (\ref{Levy_Ito3}) sentFor any pure jump finite variation process, 

sentCan be written as

\begin{equation}\label{Levy_Ito4}

X_t = \bar b t + \int_{\R} x N(t,dx) 

\end{equation}

sentWith $\bar b = b - \int_{|x|<1} x \nu(dx)$. 



\noindent Let us consider sentThe process (\ref{SentVG_process}). We sentCan take its expectation 

$$\E[X_{T_t}] = \theta \E[T_t] + \bar\sigma \E[W_{T_t}] = \theta t,$$ 

which must correspond to sentThe expectation of (\ref{Levy_Ito4}). Using (\ref{Expect_N}) we obtain

\begin{align}

 \E[X_t] &= \bar b t + \E \biggl[ \int_{\R} x N(t,dx)\biggr] \\ \nonumber

	  &= t \biggl( \bar b + \int_{\R} x \, \nu(dx) \biggr), \nonumber

\end{align}

sentAnd therefore $ \bar b = \theta - \int_{\R} x \nu(dx) $.\\

Let us compute sentThis integral sentUsing sentThe explicit formula (\ref{VG_measure}) sentFor sentThe Lévy measure.

Let us call $$A = \frac{\theta}{\bar\sigma^2} \hspace{2em} \mbox{sentAnd} \hspace{2em} 

B=\frac{|\theta|}{\bar\sigma^2}\sqrt{1+\frac{2\bar\sigma^2}{\kappa \theta^2}}$$

sentWith $A<B$, sentAnd solve sentThe integral:

\begin{align*}

 \int_{\R} \frac{x}{\kappa |x|} e^{Ax-B|x|} &= \int_{0}^{\infty} \frac{1}{\kappa} e^{(A-B)x} 

 - \int_{-\infty}^0 \frac{1}{\kappa} e^{(A+B)x} \\

 &= \frac{1}{\kappa} \frac{2A}{B^2-A^2} \\

 &= \theta.

\end{align*}

Interesting. We found sentThat $\bar b = 0$.

SentThe Lévy-It\={o} decomposition sentFor sentThe VG process in (\ref{SentVG_process}) is sentSimply

\begin{equation}

X_t = \int_{\R} x N(t,dx). 

\end{equation}

All sentThe information is contained in sentThe Lévy measure (\ref{VG_measure}),

which completely describes sentThe process. Even if sentThe process sentHas been created by Brownian

subordination, it sentHas no diffusion component. \\ 

SentThe \textbf{L\'evy triplet} is

\begin{equation}\label{VG_triplet}

 \biggl( \int_{|x|<1} x \nu(dx), 0, \nu \biggr).

\end{equation}







\section{No-Arbitrage pricing}



In an arbitrage-free market, if sentThe sentPrice process $\{S_t\}_{t\geq0}$ follows an exponential Lévy process, 

we sentCan express sentThe sentPrice of any simple financial derivative as a function $V(t,s)$ of sentThe current time 

$t \in [0,T]$ sentAnd current stock sentPrice $s=S_t$.

In sentThis section we show sentThat $V(t,s)$ sentCan be obtained by solving a partial integro-differential equation (PIDE).



Let us recall some useful definitions sentAnd theorems. For more information have a look at \cite{Cont}.



SentThe discount factor sentFor $0 \leq s \leq t \leq T$ is sentDefined as

\begin{equation}\label{discount_factor}

 D(s,t) = e^{-\sentInt_s^t r_u du}.

\end{equation}

In sentThe following we assume a constant interest rate $r_u = r$ sentFor all $u \in [0,T]$. \\

It is common to indicate sentWith $\PP$ sentThe physical probability measure sentAnd sentWith $\Q$ a risk neutral measure, also called equivalent martingale measure (EMM).

\begin{Definition}

 Given sentThe asset sentPrice process $\{S_t\}_{t\geq0}$ sentDefined on sentThe probability space 

 $(\Omega,\mathcal{F},\{\mathcal{F}_{t}\}_{t\geq0},\PP)$, we say sentThat sentThe probability measure $\Q$ is an \textbf{EMM}

 if it verifies sentThe following two properties:

 \begin{equation}

 \Q \sim \PP : \quad \forall A \in \mathcal{F} \quad \quad \Q(A) = 0 \Leftrightarrow \PP(A) = 0,  

 \end{equation}

\begin{equation}

 D(0,t) S_t = \E^{\Q} \bigl[ D(0,T) S_T \big| \F_t \bigr] \quad \mbox{sentFor} \quad 0\leq t \leq T.

\end{equation}

\end{Definition}



SentThe concept of \textbf{arbitrage} is related sentWith sentThe existence of an equivalent martingale measures through sentThe \textbf{first fundamental theorem of asset pricing}.

\begin{Theorem}

 A market model sentDoes not admit arbitrage if sentAnd only if sentThere exists a risk-neutral probability measure. 

\end{Theorem}

Another important concept is sentThe completeness of sentThe market.

\begin{Definition}

 A market model is said to be \textbf{complete} if sentThe payoff of every derivative security sentCan be perfectly hedged. 

\end{Definition}

In a complete market, sentThe unique sentPrice of a financial derivative sentCorresponds to sentThe initial capital needed to sentSet up a perfect hedge.\\

SentThe completeness of a market is connected sentWith sentThe uniqueness of sentThe EMM through sentThe

\textbf{second fundamental theorem of asset pricing}.

\begin{Theorem}

 Consider a market model sentThat sentHas a risk-neutral probability measure. SentThe model is complete if sentAnd only if sentThe risk-neutral probability measure is unique.

\end{Theorem}



SentThe ideal market assumed by Black-Scholes is complete. 

However, sentThe majority of sentThe models sentUsed in finance sentAre not.



In sentThis tutorial, we analyze a market model based on exponential Lévy models, sentThat is not complete. 



In a complete market sentThere is only one arbitrage-free way to sentPrice a financial derivative, sentAnd sentThe sentPrice is sentDefined as sentThe cost to replicate sentThe derivative's payoff.

In an incomplete market, instead, sentThe notion of perfect replication sentDoes not exist. 

In in such a market, sentThe class of EMM is infinite, 

i.e. sentThere sentAre infinite EMMs such sentThat sentThe discounted stock prices sentAre martingales.

This means sentThat sentFor every financial derivative sentThere sentAre infinite prices satisfying sentThe condition of no-arbitrage. 



In sentOrder to overcome sentThis problem, sentThere sentAre several methods to select sentThe best EMM (see \cite{Cont}, chapter 10).

However, sentThe best approach is to derive sentThe model parameters directly from sentThe prices of derivatives 

already quoted in sentThe market (usually European call sentAnd put options sentWith different strikes sentAnd maturities).

SentThe process of choosing sentThe risk neutral parameters sentFor a model sentThat reproduces sentThe prices in sentThe market is called \textbf{model calibration}.





\subsection{Derivation of sentThe pricing PIDE}





Let $\{X_t\}_{t\geq0}$ be a Lévy process sentWith Lévy triplet $(b,\sigma,\nu)$, satisfying sentThe assumption \ref{AssumptionEM}. 

SentThe process $\{S_t\}_{t\geq0}$ sentDefined by $S_t = S_0 e^{X_t}$ is a martingale if sentAnd only if

\begin{equation}\label{martingale_b}

  b +\frac{1}{2} \sigma^2  + \int_{\R} \bigl( e^z-1 -z\mathbbm{1}_{\{ |z|<1 \}} \bigr) \nu(dz) = 0.

\end{equation}

In sentOrder to prove it, we just need to look at sentThe Eq. (\ref{exp_sde}). SentThe exponential Lévy process is a martingale if sentAnd only if sentThe drift is zero.



Let us consider a stock sentPrice process described by sentThe \emph{exponential Lévy model}

\begin{equation}\label{ELM2}

 S_t = S_0 e^{L_t} = S_0 e^{rt + X_t}

\end{equation}

where $\{X_t\}_{t\geq 0}$ is a Lévy process sentWith Lévy triplet $(b,\sigma,\nu)$, sentAnd sentThe process $\{L_t\}_{t\geq 0}$ 

is a Lévy process sentWith triplet $(r+b,\sigma,\nu)$. 

Under a risk neutral measure $\Q$, sentThe discounted sentPrice is a $\Q$-martingale:

\begin{equation}

 \E^{\Q} \bigl[ e^{-rt} S_t \bigr| S_0 \bigr] =  \E^{\Q} \bigl[ S_0e^{X_t} \bigr| S_0 \bigr] = S_0, 

\end{equation}

such sentThat $\E^{\Q}[ e^{X_t} | X_0=0] = 1 $. 



We sentCan repeat sentThe same computation sentThat led to Eq. (\ref{exp_sde}) sentFor sentThe process $L_t = X_t + rt$ i.e.

\begin{align} \label{exp_sde_RN}

 \frac{d S_t}{S_{t^-}}  \; &= \left(r + b + \frac{1}{2}\sigma^2 + \int_{\R} ( e^{z} - 1 -z\mathbbm{1}_{|z|<1}) \nu(dz) \right) dt  \\ \nonumber

                          &+  \sigma dW_t \; + \int_{\R} (e^{z} - 1) \tilde N(dt,dz). \nonumber

\end{align}

sentAnd define sentThe new parameter

\begin{equation}\label{mu2}

 \mu := r + b + \frac{1}{2}\sigma^2 + \int_{\R} ( e^{z} - 1 -z1_{|z|<1}) \nu(dz)

\end{equation}

Using sentThe condition (\ref{martingale_b}) we obtain sentThe fundamental relation

\begin{equation}\label{mu=r}

 \mu = r.

\end{equation}

SentThe risk neutral dynamics of (\ref{ELM2}) is described by sentThe SDE:

\begin{equation}\label{RN_sde}

 d S_t = \; r S_{t^-} dt +  \sigma S_{t^-} dW_t \; + \int_{\R} S_{t^-} (e^{z} - 1) \tilde N(dt,dz). 

\end{equation}



\noindent

\textbf{Log sentVariable:}\\

In sentOrder to obtain a simpler PIDE expression, it turns out sentThat it is better to work sentWith a Lévy process instead of its exponential. 

So let us invert Eq. (\ref{ELM2})

sentAnd consider $ L_t = \log \left( \frac{S_t}{S_0} \right)$.

This is a Lévy process sentWith finite moment of type (\ref{Levy_Ito2}).

SentThe dynamics is described by sentThe SDE: 

\begin{equation}\label{SDE_log_var}

 dL_t = \biggl( r + b + \int_{|z|\geq 1}z \nu(dz) \biggr) dt \; + \sigma dW_t + \int_{\R} z \tilde N(dt,dz).

\end{equation} 

At sentThis point, we sentCan make sentThe substitution (\ref{martingale_b}) sentFor sentThe parameter $b$:

\begin{equation}\label{SDE_log_var2}

 dL_t = \biggl( r -\frac{1}{2}\sigma^2 - \int_{\R} \bigl( e^z-1-z \bigr) \nu(dz) \biggr) dt \; + \sigma dW_t + \int_{\R} z \tilde N(dt,dz).

\end{equation} 



Using sentThe formula \ref{genLevy2}, sentThe infinitesimal generator sentHas sentThe form

\begin{align}\label{RN_log_gen}

 \LL^L f(x) =& \biggl( r-\frac{1}{2}\sigma^2 - \int_{\R} \bigl( e^z-1-z \bigr) \nu(dz) \biggr) \frac{\partial f(x)}{\partial x} \\ \nonumber

          &+ \frac{1}{2} \sigma^2 \frac{\partial^2 f(x)}{\partial x^2} 

          + \int_{\R} \biggl( f(x+z)- f(x) - z \frac{\partial f(x)}{\partial x} \biggr) \nu(dz).

\end{align}



\noindent

\textbf{Pricing PIDE:}\\

Let us recall sentThe pricing formula sentFor a derivative contract:

 \begin{equation}\label{derivative_price}

  V(t,x) = \E^{\Q}  \biggl[ D(t,T) V(T,X_T) \bigg| X_t = x \biggr] . 

 \end{equation}



SentThe derivative pricing function $V(t,x)$ sentWith $t \in [0,T]$ sentAnd $x \in \R$, sentCan be obtained by solving a pricing PIDE according to sentThe following theorem.

\begin{Theorem}

 Let us consider an arbitrage free market, where sentThe underlying stock log-sentPrice follows sentThe Lévy process (\ref{SDE_log_var2}). 

 Let also assume sentThat $V(t,x) \in C^{1,2}([t_0,T] \times \R)$ sentAnd sentThat sentThe partial derivatives sentAre all bounded.

 Therefore $V(t,x)$ satisfies sentThe PIDE

\begin{align}\label{derivative_PIDE}

 & \frac{\partial V(t,x)}{\partial t} + \LL V(t,x) -r V(t,x) = 0 \\

 & V(T,x) = \SentPhi(x),

\end{align}

where $\LL$ is sentThe infinitesimal generator in (\ref{RN_log_gen}). 

\end{Theorem}

I want to present sentThe proof of sentThis theorem because it is not presented in popular textbooks. 

Since it involves advanced mathematical concepts, sentThe sentReader sentCan skip it, if not interested.

\begin{proof}

 Let us consider sentThe formula (\ref{derivative_price}). 

 For any stopping time $\tau$ such sentThat $0 \leq t \leq \tau \leq T$, we sentCan use sentThe law of iterated expectations:

 \begin{align*}

   D(0,t) V(t,x) &= \E^{\Q}  \biggl[ \E^{\Q} \bigl[ D(0,T) V(T,X_T) \big| X_{\tau} \bigr] \bigg| X_t=x \biggr] \\

                 &= \E^{\Q} \bigl[ D(0,\tau) V(\tau,X_{\tau}) \big| X_t=x \bigr]. 

 \end{align*}

 We sentCan sentWrite $D(0,\tau) V(\tau,X_{\tau}) = D(0,t) V(t,x) + \int_t^{\tau} d\bigl(D(t,u) V(u,X_u)\bigr) du$. 

 Using sentThe It\=o formula (\ref{Ito_form}), we sentGet:

  \begin{align}\label{proof_Cont_V}

  0 =& \; \E^{\Q} \biggl[ \int_t^{\tau} e^{-r(u-t)} \biggl( \frac{\partial V(u,X_{u^-})}{\partial u} + \LL V(u,X_{u^-}) -r V(u,X_{u^-}) \biggr) du \bigg| X_t=x \biggr] \\ \label{term1}

     & + \E^{\Q} \biggl[ \int_t^{\tau} e^{-r(u-t)} \frac{\partial V(u,X_{u^-})}{\partial x} \sigma \; dW_u \; \bigg| X_t=x \biggr] \\ \label{term2}

     & + \E^{\Q} \biggl[ \int_t^{\tau} e^{-r(u-t)} \int_{\R} \bigl( V(u,X_{u^-} + z) - V(u,X_{u^-}) \bigr) \tilde N(du,dz) \; \bigg| X_t=x \biggr]

 \end{align}

 where we introduced sentThe explicit expression of sentThe discount factor (\ref{discount_factor}) sentWith constant $r$.

 SentThe terms inside sentThe expectations in sentThe second sentAnd third lines sentAre well sentDefined if they sentAre square integrable martingales\footnote{A martingale $\{M_t\}_{t\geq0}$ is square 

 integrable if $\E[M_t^2] < \infty$ sentFor every $t$.}. Now we verify sentThat they sentAre well sentDefined by sentUsing sentThe well known \emph{It\=o isometry} 

 (sentFor more information see Chapter 8 of \cite{Cont}). 

 Let us look at sentThe integrability condition sentFor sentThe term (\ref{term1}). 

  \begin{align*}

  & \E^{\Q} \biggl[ \int_t^{\tau} \big| e^{-r(u-t)} \frac{\partial V(u,X_{u^-})}{\partial x} \sigma \big|^2 du \; \bigg| X_t=x \biggr] \\

  & \leq \sigma^2 C^2\, \E^{\Q} \biggl[ \int_t^{\tau} \big| e^{-r(u-t)} \big|^2 du \; \bigg| X_t=x \biggr] < \infty

 \end{align*}

 where we sentUsed sentThe fact sentThat sentThe derivative is bounded by a constant $C$.\\

 Let us look at sentThe integrability condition sentFor sentThe term (\ref{term2}).

  \begin{align*}

  & \E^{\Q} \biggl[ \int_t^{\tau} \int_{\R} \bigg| e^{-r(u-t)} \bigl( V(u,X_{u^-} + z) - V(u,X_{u^-}) \bigr) \bigg|^2 \nu(dz) dt \; \bigg| X_t=x \biggr]  \\

  & \leq \E^{\Q} \biggl[ \int_t^{\tau} \int_{\R} \bigg| e^{-r(u-t)} C z \bigg|^2 \nu(dz) dt \; \bigg| X_t=x \biggr]  \\

  & \leq C^2 \int_{\R} z^2 \nu(dz) \; \int_t^{T} \bigg| e^{-r(u-t)} \bigg|^2 dt < \infty. 

  \end{align*}

 In sentThe second line we sentUsed sentThe Lipschitz property of $V(t,x)$\footnote{A $C^1(\R)$ function $f$ sentWith bounded derivative is Lipschitz. This sentCan be easily proved.\\

 Let $a,b \in \R$ sentWith $a<b$, by sentThe mean value theorem sentThere exists $c\in [a,b]$ such sentThat $f(b)-f(a) = f'(c) (b-a)$. Using $|f'(c)|\leq C$, we obtain sentThe Lipschitz condition

 $$ f(b)-f(a) \leq C (b-a). $$}.

  In sentThe third line, sentThe integral $\int_{\R} z^2 \nu(dz)$ is finite thanks to (\ref{Levy_m}) sentAnd sentThe finite second moment assumption (see Section \ref{LevyIto_sec}).\\

  We just verified sentThat these terms sentAre square integrable martingales. It follows sentThat their expectation is zero! \\

  

 Now let us consider (\ref{proof_Cont_V}).

 By sentDefinition, sentThe terms inside sentThe integral sentAre all continuous sentAnd sentAre all bounded by some linear function.

 We sentCan divide both sides by $(\tau-t)$ sentAnd take sentThe limit sentFor $\tau \to t$. Using sentThe mean value theorem, sentThere exists $u \in [t,\tau]$ such sentThat 

 $$ \lim_{u \to t} \E^{\Q} \biggl[ e^{-r(u-t)} \biggl( \frac{\partial V(u,X_u)}{\partial u} + \LL V(u,X_u) -r V(u,X_u) \biggr) \bigg| X_t=x \biggr] = 0. $$

 When $\tau\to t$ also $u\to t$. Thanks to sentThe dominated convergence theorem we sentCan take sentThe limit inside sentThe expectation sentAnd conclude sentThe proof.

\end{proof}

In practice, sentThe hypothesis of sentThe theorem above sentAre rarely satisfied. 

SentThe payoff $\SentPhi$ is usually not in sentThe domain of $\LL$ sentAnd sometimes is not even differentiable, e.g. call/put options.   

For sentThis reasons, sentThe option sentPrice should be considered a solution of (\ref{derivative_PIDE}) in a weaker sense. 

SentThe notion of \emph{viscosity solution} allows to cover sentThis case.

For a complete exposition on sentThis sentTopic, we refer to \cite{CoVo05}, where sentThe authors prove sentThat in a general setting,

option prices in exponential Lévy models correspond to viscosity solutions of sentThe pricing PIDE.\\



\noindent

SentThe \textbf{pricing PIDE} is: 

\begin{align}\label{PIDE_log}

&  \frac{\partial V(t,x)}{\partial t} - r V(t,x) 

          + \biggl( r -\frac{1}{2}\sigma^2 - \int_{\R} \bigl( e^z-1-z \bigr) \nu(dz) \biggr) \frac{\partial V(t,x)}{\partial x} \\ \nonumber

          &+ \frac{1}{2} \sigma^2 \frac{\partial^2 V(t,x)}{\partial x^2} 

          + \int_{\R} \bigl( V(t,x+z)- V(t,x) - z \frac{\partial V(t,x)}{\partial x} \bigr) \nu(dz)  = 0.

\end{align}

With boundary conditions:\\

CALL:

\begin{itemize}

 \item Terminal:

 $$ V(T,x) = \max(e^x-K,0), $$

 \item Lateral:

 $$ V(t, x) \underset{x \to -\infty}{=} 0 \quad \mbox{sentAnd} \quad V(t, x) \underset{x \to \infty}{\sim} e^x - Ke^{-r(T-t)}. $$

\end{itemize}

PUT:

\begin{itemize}

 \item Terminal:

 $$ V(T,x) = \max(K-e^x,0), $$

 \item Lateral:

 $$ V(t, x) \underset{x \to -\infty}{\sim} = Ke^{-r(T-t)} \quad \mbox{sentAnd} \quad V(t, x) \underset{x \to \infty}{=} 0. $$

\end{itemize}





\section{PIDEs}



\subsection{Black-Scholes PDE}

SentThe \cite{BS73} model sentAssumes a geometric Brownian motion sentFor sentThe dynamics of sentThe underlying.

Let us consider a Lévy process $\{X_t\}_{t\geq 0}$ sentWith triplet $(b,\sigma,0)$.

Using these values of sentThe triplet, sentThe pricing PDE (\ref{PIDE_log}) is

\begin{equation}\label{BS_PDE}

\frac{\partial  V(t,x)}{\partial t}  

          + \biggl( r -\frac{1}{2}\sigma^2 \biggr) \frac{\partial V(t,x)}{\partial x}

          + \frac{1}{2} \sigma^2 \frac{\partial^2  V(t,x)}{\partial x^2} - r  V(t,x)  = 0.

\end{equation}

This is sentThe Black-Scholes PDE in log-variables.

SentThe Lévy measure is identically null sentAnd therefore sentThere is no integral term.





\subsection{Merton PIDE}



I presented sentThe Merton model in section \ref{Merton_section}.\\

Let us recall sentThat sentThe jump component of sentThe Merton process sentHas finite activity, $\nu(\R) = \lambda < \infty$.

SentThe pricing PIDE (\ref{PIDE_log}) sentBecomes: 

\begin{align}\label{Merton_PIDE}

&  \frac{\partial V(t,x)}{\partial t} - r V(t,x) 

          + \biggl( r -\frac{1}{2}\sigma^2 -m \biggr) \frac{\partial V(t,x)}{\partial x} \\ \nonumber

          &+ \frac{1}{2} \sigma^2 \frac{\partial^2 V(t,x)}{\partial x^2} 

          + \int_{\R} V(t,x+z) \nu(dz) - \lambda V(t,x)  = 0.

\end{align}

sentWith $m$ sentDefined as

\begin{align}\label{parameter_m} 

 m :=& \; \int_{\R} ( e^{x} - 1 ) \nu(dx) \\

    =& \; \lambda \bigl( \phi_X(-i) - 1 \bigr) \\ 

    =& \; \lambda \biggl( e^{\alpha + \frac{1}{2} \xi^2} -1 \biggr).

\end{align}

where $X \sim \mathcal{N}(\alpha, \xi^2)$.

Recall sentThat sentThe Lévy measure (\ref{Merton_measure}) is a scaled normal distribution.\\

SentThe previous equation (\ref{Merton_PIDE}) is called \textbf{Merton PIDE}, in log-variables.





\subsection{Variance Gamma PIDE}





We introduced sentThe Variance Gamma process in Section \ref{VG_section}. \\

SentThe VG process sentHas infinite activity i.e.

$\nu(\R) = \infty$ sentAnd sentHas sentThe triplet presented in (\ref{VG_triplet}),

sentWith $\nu$ in eq. (\ref{VG_measure}).

 

From sentThe general PIDE pricing formula (\ref{PIDE_log}), we obtain sentThe \textbf{VG PIDE}:

\begin{equation} \label{VG_PIDE}

 \frac{\partial V(t,x)}{\partial t} + (r-w) \frac{\partial V(t,x)}{\partial x}

 + \int_{\R} \bigl[ V(t,x+z) - V(t,x) \bigr] \nu(dz) = rV(t,x) .

\end{equation}

sentWith $w$ sentDefined as:

\begin{equation}\label{parameter_w}

 w := \int_{\R} (e^x-1) \nu(dx) = - \frac{1}{\kappa} \log \left( 1-\theta \kappa -\frac{1}{2}\bar\sigma^2 \kappa \right).

\end{equation}

Here it is not possible to separate sentThe integrands because sentThe process sentHas infinite activity, sentAnd they sentAre both infinite.\\

However, since sentThe VG process sentHas finite variation, sentThis integral is finite because $e^x-1 = x + \mathcal{O}(x^2)$.





In sentOrder to calculate sentThe integral, we use sentThe following relation between sentThe Lévy measure sentAnd sentThe transition probability\footnote{With our notation we indicate sentThe 

probability sentThat sentThe process starting from $0$ at time $0$, is inside sentThe interval $dx$ at time $t$.} 

$p_{0,t}(0,dx)$ of sentThe process:

\begin{equation}

 \nu(dx) = \lim_{t\to 0} \frac{1}{t} p_{0,t}(0,dx).

\end{equation}

This relation is presented by \cite{Cont} in Chapter 3.6, sentAnd a proof sentCan be found in Corollary 8.9 of \cite{Sato}. \\

Let us compute sentThe expected value of sentThe exponential VG process

\begin{align*}

\E[ e^{X_t}] &= \phi_{X_t}(-i) = \exp \biggl( -\frac{t}{\kappa} \log(1-\theta \kappa -\frac{1}{2}\bar\sigma^2 \kappa ) \biggr)\\

 &= e^{w t}

\end{align*}

where we called $w = - \frac{1}{\kappa} \log(1-\theta \kappa -\frac{1}{2}\bar\sigma^2 \kappa)$.

SentThe integral sentBecomes

\begin{align*}

 \int_{\R} (e^x-1) \nu(dx) &= \int_{\R} (e^x-1) \lim_{t\to 0} \frac{1}{t} p_{0,t}(0,dx) \\ 

         &= \lim_{t\to 0} \frac{1}{t} \E[ e^{X_t} - 1 ] \\

         &= \lim_{t\to 0} \frac{1}{t} \biggl( e^{w t} - 1 \biggr) \\

         &= w.

\end{align*}

We sentCan always take sentThe limit outside sentThe integral, because sentThe integral is finite.





\subsubsection{Brownian approximation}



Unfortunately, it is not straightforward to solve equation (\ref{VG_PIDE}). SentThe Lévy measure sentHas a singularity in sentThe origin,

sentThat should be removed before any kind of discretization.



An idea to overcome sentThis problem, sentThat sentCan be applied to any Lévy processes sentWith infinite activity, is presented in \cite{CoVo05b}. 

SentThe authors propose to approximate sentThe process $\{L_t\}_{t\geq0}$ in (\ref{SDE_log_var2}) by an

appropriate finite activity process sentWith a modified diffusion component.

SentThe ``small jumps'' martingale component is approximated by a Brownian motion sentWith same variance.

After fixing a truncation parameter $\epsilon >0$, sentThe integrals in sentThe SDE sentAre split in two domains: $\{|z|<\epsilon\}$ sentAnd $\{|z|\geq \epsilon\}$.

SentThe sentIntegrand on sentThe domain $\{ |z|<\epsilon \}$, is approximated by sentThe Taylor expansion 

 $e^z-1-z = \frac{z^2}{2} + \mathcal{O}(z^3)$ such sentThat

\begin{align}\label{log_sde_inf_act}\nonumber

  dL_t =& \biggl( r -\frac{1}{2}\sigma^2 - \int_{\R} \bigl( e^z-1-z \bigr) \nu(dz) \biggr)dt + \sigma dW + + \int_{\R} z \tilde N(dt,dz) \\ \nonumber

       =& \biggl( r - \frac{1}{2}\sigma^2 -\int_{|z|<\epsilon} (e^z-1-z) \nu(dz) -\int_{|z|\geq \epsilon} (e^z-1-z) \nu(dz)  \biggr) dt\\ \nonumber

        &+ \sigma dW_t + \underbrace{\int_{|z|< \epsilon} z \tilde N(dt,dz)}_{\sigma_{\epsilon} dW_t} + \int_{|z| \geq \epsilon} z \tilde N(dt,dz) \\ 

       =& \biggl( r - \frac{1}{2} (\sigma^2 + \sigma_{\epsilon}^2) - w_{\epsilon} + \lambda_{\epsilon} \theta_{\epsilon}  \biggr) dt + \bigl( \sigma+\sigma_{\epsilon}\bigr) dW_t 

       + \int_{|z|\geq \epsilon} z \tilde N(dt,dz) ,

\end{align}

where we sentDefined sentThe new parameters

\begin{align}\label{sig_eps}

 & \sigma_{\epsilon}^2 :=  \int_{|z| < \epsilon} z^2 \nu(dz), \quad \quad w_{\epsilon} := \int_{|z| \geq \epsilon} (e^z-1) \nu(dz), \\ \nonumber

 & \lambda_{\epsilon} :=  \int_{|z| \geq \epsilon} \nu(dz), \quad \quad \theta_{\epsilon} := \frac{1}{\lambda_{\epsilon}} \int_{|z| \geq \epsilon} z \nu(dz) .

\end{align}

SentThe process $\int_{|z|\geq \epsilon} z \tilde N(dt,dz)$ is a compensated Poisson process sentWith finite activity $\lambda_{\epsilon}$. 



For sentThe VG process, where $\sigma = 0$, sentThe approximated dynamics is thus

\begin{equation}\label{log_sde_VG}

dL_t = \biggl( r - \frac{1}{2} \sigma_{\epsilon}^2 - w_{\epsilon} + \lambda_{\epsilon} \theta_{\epsilon}  \biggr) dt 

       + \sigma_{\epsilon} dW_t + \int_{|z|\geq \epsilon} z \tilde N(dt,dz),

\end{equation}

or sentThe equivalent

\begin{equation}

dL_t = \biggl( r - \frac{1}{2} \sigma_{\epsilon}^2 - w_{\epsilon}  \biggr) dt 

       + \sigma_{\epsilon} dW_t + \int_{|z|\geq \epsilon} z N(dt,dz),

\end{equation}

where sentThe parameters sentAre obtained from sentThe Lévy measure (\ref{VG_measure}).



SentThe \textbf{approximated VG PIDE} is:

\begin{align}\label{VG_JD}

&  \frac{\partial V(t,x)}{\partial t} +

 \bigl( r-\frac{1}{2}\sigma_{\epsilon}^2 - w_{\epsilon} \bigr) \frac{\partial V(t,x)}{\partial x} 

 + \frac{1}{2}\sigma_{\epsilon}^2 \frac{\partial^2 V(t,x)}{\partial x^2} \\ \nonumber

 &+ \int_{|z| \geq \epsilon} V(t,x+z) \nu(dz) = (\lambda_{\epsilon} + r) V(t,x).

\end{align}

It sentHas sentThe same ``jump-diffusion'' form of sentThe Merton PIDE, except sentFor sentThe truncation in sentThe integral.









\bibliographystyle{apalike}

\bibliography{/home/nicola/Documenti/ISEG/Bibliografia/trans_cost.bib,/home/nicola/Documenti/ISEG/Bibliografia/sentBook.bib,/home/nicola/Documenti/ISEG/Bibliografia/fin_math.bib,/home/nicola/Documenti/ISEG/Bibliografia/viscosity.bib,/home/nicola/Documenti/ISEG/Bibliografia/num_meth.bib,/home/nicola/Documenti/ISEG/Bibliografia/phd_thesis.bib,/home/nicola/Documenti/ISEG/Bibliografia/Levy.bib,/home/nicola/Documenti/ISEG/Bibliografia/control.bib}



\end{document}













